% Options for packages loaded elsewhere
\PassOptionsToPackage{unicode}{hyperref}
\PassOptionsToPackage{hyphens}{url}
\PassOptionsToPackage{dvipsnames,svgnames,x11names}{xcolor}
%
\documentclass[
  10pt,
]{article}
\usepackage{amsmath,amssymb}
\usepackage{iftex}
\ifPDFTeX
  \usepackage[T1]{fontenc}
  \usepackage[utf8]{inputenc}
  \usepackage{textcomp} % provide euro and other symbols
\else % if luatex or xetex
  \usepackage{unicode-math} % this also loads fontspec
  \defaultfontfeatures{Scale=MatchLowercase}
  \defaultfontfeatures[\rmfamily]{Ligatures=TeX,Scale=1}
\fi
\usepackage{lmodern}
\ifPDFTeX\else
  % xetex/luatex font selection
\fi
% Use upquote if available, for straight quotes in verbatim environments
\IfFileExists{upquote.sty}{\usepackage{upquote}}{}
\IfFileExists{microtype.sty}{% use microtype if available
  \usepackage[]{microtype}
  \UseMicrotypeSet[protrusion]{basicmath} % disable protrusion for tt fonts
}{}
\makeatletter
\@ifundefined{KOMAClassName}{% if non-KOMA class
  \IfFileExists{parskip.sty}{%
    \usepackage{parskip}
  }{% else
    \setlength{\parindent}{0pt}
    \setlength{\parskip}{6pt plus 2pt minus 1pt}}
}{% if KOMA class
  \KOMAoptions{parskip=half}}
\makeatother
\usepackage{xcolor}
\usepackage[margin=1in]{geometry}
\usepackage{color}
\usepackage{fancyvrb}
\newcommand{\VerbBar}{|}
\newcommand{\VERB}{\Verb[commandchars=\\\{\}]}
\DefineVerbatimEnvironment{Highlighting}{Verbatim}{commandchars=\\\{\}}
% Add ',fontsize=\small' for more characters per line
\newenvironment{Shaded}{}{}
\newcommand{\AlertTok}[1]{\textcolor[rgb]{1.00,0.00,0.00}{\textbf{#1}}}
\newcommand{\AnnotationTok}[1]{\textcolor[rgb]{0.38,0.63,0.69}{\textbf{\textit{#1}}}}
\newcommand{\AttributeTok}[1]{\textcolor[rgb]{0.49,0.56,0.16}{#1}}
\newcommand{\BaseNTok}[1]{\textcolor[rgb]{0.25,0.63,0.44}{#1}}
\newcommand{\BuiltInTok}[1]{\textcolor[rgb]{0.00,0.50,0.00}{#1}}
\newcommand{\CharTok}[1]{\textcolor[rgb]{0.25,0.44,0.63}{#1}}
\newcommand{\CommentTok}[1]{\textcolor[rgb]{0.38,0.63,0.69}{\textit{#1}}}
\newcommand{\CommentVarTok}[1]{\textcolor[rgb]{0.38,0.63,0.69}{\textbf{\textit{#1}}}}
\newcommand{\ConstantTok}[1]{\textcolor[rgb]{0.53,0.00,0.00}{#1}}
\newcommand{\ControlFlowTok}[1]{\textcolor[rgb]{0.00,0.44,0.13}{\textbf{#1}}}
\newcommand{\DataTypeTok}[1]{\textcolor[rgb]{0.56,0.13,0.00}{#1}}
\newcommand{\DecValTok}[1]{\textcolor[rgb]{0.25,0.63,0.44}{#1}}
\newcommand{\DocumentationTok}[1]{\textcolor[rgb]{0.73,0.13,0.13}{\textit{#1}}}
\newcommand{\ErrorTok}[1]{\textcolor[rgb]{1.00,0.00,0.00}{\textbf{#1}}}
\newcommand{\ExtensionTok}[1]{#1}
\newcommand{\FloatTok}[1]{\textcolor[rgb]{0.25,0.63,0.44}{#1}}
\newcommand{\FunctionTok}[1]{\textcolor[rgb]{0.02,0.16,0.49}{#1}}
\newcommand{\ImportTok}[1]{\textcolor[rgb]{0.00,0.50,0.00}{\textbf{#1}}}
\newcommand{\InformationTok}[1]{\textcolor[rgb]{0.38,0.63,0.69}{\textbf{\textit{#1}}}}
\newcommand{\KeywordTok}[1]{\textcolor[rgb]{0.00,0.44,0.13}{\textbf{#1}}}
\newcommand{\NormalTok}[1]{#1}
\newcommand{\OperatorTok}[1]{\textcolor[rgb]{0.40,0.40,0.40}{#1}}
\newcommand{\OtherTok}[1]{\textcolor[rgb]{0.00,0.44,0.13}{#1}}
\newcommand{\PreprocessorTok}[1]{\textcolor[rgb]{0.74,0.48,0.00}{#1}}
\newcommand{\RegionMarkerTok}[1]{#1}
\newcommand{\SpecialCharTok}[1]{\textcolor[rgb]{0.25,0.44,0.63}{#1}}
\newcommand{\SpecialStringTok}[1]{\textcolor[rgb]{0.73,0.40,0.53}{#1}}
\newcommand{\StringTok}[1]{\textcolor[rgb]{0.25,0.44,0.63}{#1}}
\newcommand{\VariableTok}[1]{\textcolor[rgb]{0.10,0.09,0.49}{#1}}
\newcommand{\VerbatimStringTok}[1]{\textcolor[rgb]{0.25,0.44,0.63}{#1}}
\newcommand{\WarningTok}[1]{\textcolor[rgb]{0.38,0.63,0.69}{\textbf{\textit{#1}}}}
\setlength{\emergencystretch}{3em} % prevent overfull lines
\providecommand{\tightlist}{%
  \setlength{\itemsep}{0pt}\setlength{\parskip}{0pt}}
\setcounter{secnumdepth}{-\maxdimen} % remove section numbering
\ifLuaTeX
  \usepackage{selnolig}  % disable illegal ligatures
\fi
\usepackage{bookmark}
\IfFileExists{xurl.sty}{\usepackage{xurl}}{} % add URL line breaks if available
\urlstyle{same}
\hypersetup{
  colorlinks=true,
  linkcolor={black},
  filecolor={Maroon},
  citecolor={Blue},
  urlcolor={black},
  pdfcreator={LaTeX via pandoc}}

\author{}
\date{}

\begin{document}

\section{DDTA research work plan after documentation-layer
closure}\label{ddta-research-work-plan-after-documentation-layer-closure}

\textbf{WORK PLAN - REVISION 5}

\textbf{Status:} active after BA0-T3 closure.

\textbf{Supersedes:} Revision 4 only for forward execution state; R1-R4
remain historical research records.

\subsection{1. Fixed prior state}\label{fixed-prior-state}

\begin{itemize}
\tightlist
\item
  Chapters 2-4: CLOSED / FINAL for current thesis scope.
\item
  Documentation layer: CLOSED.
\item
  BA0-R systems-modeling prior-art research: CLOSED.
\item
  BA0-T1: COMPLETED / PROVISIONAL.
\item
  BA0-T2: COMPLETED / PROVISIONAL PASS.
\item
  BA0-T3: CLOSED.
\item
  BA0 responsibility and non-goals: CLOSED.
\item
  No BAE type or relation was accepted by BA0.
\item
  ThreatForge remains a case study/tooling instrument, not DDTA semantic
  authority.
\end{itemize}

\subsection{2. Closed Base Analysis
responsibility}\label{closed-base-analysis-responsibility}

The authoritative BA0 responsibility boundary for forward methodology
work is \href{BA0_BASE_ANALYSIS_RESPONSIBILITY_BOUNDARY_R1.md}{Base
Analysis responsibility boundary R1}.

Base Analysis is the accepted methodology-neutral analytical
representation of shared project meaning required by DDTA. It preserves
source/origin provenance, baseline-scoped semantic identity and explicit
unresolved state sufficiently to support shared human/method
projections, source drill-down, change-impact reasoning and
source-localized feedback handoff.

Governed documentation remains project authority. Method-specific
interpretation, analysis results, corrective policy decisions and
documentation governance remain outside the shared Base Analysis core.

Semantic canonicality does not require one permanently materialized
graph.

\subsection{3. Closed BA0 responsibility
set}\label{closed-ba0-responsibility-set}

Forward work must preserve:

\begin{enumerate}
\def\labelenumi{\arabic{enumi}.}
\tightlist
\item
  BA0-C1 Authority and provenance boundary.
\item
  BA0-C2 Baseline-scoped shared semantic identity.
\item
  BA0-C3 Method-neutral shared core.
\item
  BA0-C4 Explicit uncertainty and diagnostic localization.
\item
  BA0-C5 Projection readiness and source drill-down.
\item
  BA0-C6 Change-impact traceability.
\item
  BA0-C7 Source-localized feedback handoff.
\item
  BA0-C8 Minimality and representation independence.
\end{enumerate}

\texttt{grounded}, \texttt{derived} and \texttt{diagnostic/unresolved}
are required responsibility-level distinctions. A generic
\texttt{reviewed\ analytical\ addition} class is not required by BA0 and
must not be assumed by BA1.

\subsection{4. Sequence after BA0
closure}\label{sequence-after-ba0-closure}

Do not skip phases:

\begin{Shaded}
\begin{Highlighting}[]
\NormalTok{BA1 {-} Minimal BAE ontology}
\NormalTok{BA2 {-} Relations and canonical action vocabulary}
\NormalTok{BA3 {-} Document{-}to{-}BAE derivation, provenance and authority}
\NormalTok{BA4 {-} Multi{-}level human and method projections}
\NormalTok{BA5 {-} Controlled vocabulary and optional authoring/extraction assistance}
\NormalTok{BA6 {-} Base Analysis closure and regression}
\end{Highlighting}
\end{Shaded}

BA0 closure authorizes ontology work; it does not pre-approve any type.

\subsection{5. BA1 - Minimal BAE
ontology}\label{ba1---minimal-bae-ontology}

\subsubsection{BA1-T1 - minimal BAE ontology candidate
derivation}\label{ba1-t1---minimal-bae-ontology-candidate-derivation}

\textbf{NEXT / NOT EXECUTED BY THIS PLAN REVISION.}

Start from BA0-C1..C8, BA0-R prior-art pressure and governed DDTA
corpora. For each recurring semantic responsibility required by the
closed boundary, test whether it needs:

\begin{itemize}
\tightlist
\item
  an independently identified entity;
\item
  a relation;
\item
  a role;
\item
  a property/attribute;
\item
  a constraint/state;
\item
  a provenance/diagnostic record;
\item
  or only a derived/projection construct.
\end{itemize}

The first pass must prefer the smallest representation that preserves
the closed responsibilities. \texttt{semantic\ responsibility\ exists}
does not imply \texttt{first-class\ metaclass}.

Historical \texttt{actor/component/asset/boundary/data\_flow}
categories, STRIDE DFD constructs, ThreatForge implementation classes
and SysML/KerML/ArchiMate/AADL concepts are candidate comparisons only.
None is accepted by inheritance.

\subsubsection{BA1-T1 exit}\label{ba1-t1-exit}

Produce a falsifiable minimal ontology candidate with explicit evidence
for every proposed first-class identity and an explicit
rejected/deferred list. Do not begin BA2 relation-vocabulary closure
until the BA1 candidate has been pressure-tested.

\subsection{6. BA2 - Relations and canonical action
vocabulary}\label{ba2---relations-and-canonical-action-vocabulary}

After BA1, define the smallest stable semantic relation vocabulary. Keep
semantic relation, canonical lexical label and authoring synonym
assistance separate.

\subsection{7. BA3 - Document-to-BAE derivation, provenance and
authority}\label{ba3---document-to-bae-derivation-provenance-and-authority}

Close exact identity/equivalence/lifecycle rules, source locators,
grounded/derived/diagnostic material states, review boundaries, stale
state and document-to-analysis traceability.

If a methodology-neutral shared-core analytical addition becomes
demonstrably necessary here, do not silently introduce it. Reassess
whether BA0 must reopen.

Automatic candidate extraction may be evaluated only as optional
assistance and is not required for thesis validity.

\subsection{8. BA4 - Multi-level human and method
projections}\label{ba4---multi-level-human-and-method-projections}

Define projections over the same accepted Base Analysis semantics. Test
progressive human navigation and method-specific views without creating
a second authored project model. No view is source authority.

\subsection{9. BA5 - Controlled vocabulary and optional
assistance}\label{ba5---controlled-vocabulary-and-optional-assistance}

Vocabulary reuse, duplicate suggestions, NLP/LLM proposals and candidate
extraction are optional support. None may silently establish identity,
equivalence, ownership, project fact or BAE acceptance.

\subsection{10. BA6 - Base Analysis closure and
regression}\label{ba6---base-analysis-closure-and-regression}

Regress closed corpora and at least one structurally different holdout.
Verify semantic identity stability, source traceability, projection
consistency, diagnostic behavior, change/re-analysis scope, method
neutrality and human drill-down support.

BA6 closes the implemented Base Analysis design; BA0 has already closed
only the responsibility boundary.

\subsection{11. Analysis envelope after Base Analysis
closure}\label{analysis-envelope-after-base-analysis-closure}

Only after BA6:

\begin{itemize}
\tightlist
\item
  A1 - AnalysisRecord / execution envelope;
\item
  A2 - common reviewed Finding boundary;
\item
  A3 - change/provenance integration.
\end{itemize}

Then define the generic overlay contract and full methodology
evaluations.

\subsection{12. Method-specific
evaluations}\label{method-specific-evaluations}

Planned concrete evaluations remain STRIDE and STRIDE-AI over the same
common core. They are demonstrators of the overlay boundary, not proof
of universal method support.

The BA0-T2 STRIDE probe remains only a pre-overlay consumer probe and
does not count as formal STRIDE evaluation.

\subsection{13. Corrective feedback
boundary}\label{corrective-feedback-boundary}

The general DDTA loop remains:

\begin{Shaded}
\begin{Highlighting}[]
\NormalTok{analysis / validation / diagnostic}
\NormalTok{ {-}\textgreater{} source{-}localized corrective candidate}
\NormalTok{ {-}\textgreater{} governed review}
\NormalTok{ {-}\textgreater{} correction / clarification / supersession / new requirement or decision}
\NormalTok{ {-}\textgreater{} Base Analysis revalidation}
\NormalTok{ {-}\textgreater{} impacted re{-}analysis}
\end{Highlighting}
\end{Shaded}

BA0-C7 requires traceable handoff to the governed source area. It does
not assign the review/acceptance workflow to Base Analysis.

The thesis may still evaluate the narrower accepted
security-Finding-to-SecurityRequirement path as a specific research
contribution.

\subsection{14. Next authorized
microstep}\label{next-authorized-microstep}

Execute only:

\begin{quote}
\textbf{BA1-T1 - minimal BAE ontology candidate derivation from the
closed BA0 responsibilities.}
\end{quote}

Do not start BA2, formal STRIDE overlay design, Common Finding schema or
implementation work during this microstep.

\end{document}
